同一份代码，同一份数据，同一组超参数，只把随机种子从 \(0\) 换成 \(1\)，测试准确率从 \(76.2\%\) 变成 \(74.8\%\)。

没有 bug。这个差距比上周那次``改进''带来的提升还大。

这件事在前面十三个部分里几乎不该发生。那里的每个结论都附着一组条件：条件满足则结论成立，条件不满足则给出反例。到了这一部分，最诚实的回答常常变成``实践中这样做效果更好''——于是我们第一次要面对这样一类知识：它是真的，它有用，但它不是定理。

这不是这一部分写得不够严谨，是它研究的对象本身处在这个状态。深度学习真正的结构性变化只有一条：\textbf{它把``用什么表示''从人工设计变成了可优化的对象。}前面那些模型族里，特征由人挑好，模型只负责在给定的特征上拟合，于是``假设了什么''是写在纸面上的；深度模型让表示本身成为参数，假设随之沉进权重里，再也读不出来。层数、激活函数、优化器的种种选择，都是这一条的推论；而所有关于``为什么有效''的保证退化成经验规律，是这一条的代价。

\begin{center}
\begin{tikzpicture}[x=1cm,y=1cm,font=\small,>=stealth]
  \draw[->,black!60,line width=0.9pt] (0,0) -- (13,0);
  \node[font=\footnotesize,black!70,anchor=north west] at (0,-0.15) {有定理};
  \node[font=\footnotesize,black!70,anchor=north east] at (13,-0.15) {只有反复观察};

  \foreach \x in {0.9,3.1,5.6,8.3,10.6,12.2} \draw[black!45] (\x,0) -- (\x,0.28);

  \node[align=center,font=\footnotesize,black!85,fill=white] at (0.9,0.72)
    {反向传播\\＝链式法则};
  \node[align=center,font=\footnotesize,black!85,fill=white] at (3.1,1.62)
    {卷积的平移等变\\注意力的置换等变};
  \node[align=center,font=\footnotesize,black!85,fill=white] at (5.6,0.72)
    {通用逼近定理\\（存在，不构造）};
  \node[align=center,font=\footnotesize,black!85,fill=white] at (8.3,1.62)
    {残差保住梯度通路\\（有分析，有经验）};
  \node[align=center,font=\footnotesize,black!85,fill=white] at (10.6,0.72)
    {缩放律\\（拟合出的幂律）};
  \node[align=center,font=\footnotesize,black!85,fill=white] at (12.2,1.62)
    {过参数化\\为何仍能泛化};

  \draw[black!35] (0.9,0.28) -- (0.9,0.42);
  \draw[black!35] (3.1,0.28) -- (3.1,1.2);
  \draw[black!35] (5.6,0.28) -- (5.6,0.42);
  \draw[black!35] (8.3,0.28) -- (8.3,1.2);
  \draw[black!35] (10.6,0.28) -- (10.6,0.42);
  \draw[black!35] (12.2,0.28) -- (12.2,1.2);
\end{tikzpicture}

\emph{同一部分里的结论，认识论地位相差很远。左端的东西可以证明，右端的东西只是被反复观察到。把两者混为一谈——尤其是把右端的规律当成定理往外推——是这个领域里最贵的一类判断失误。}
\end{center}

\subsection{五条贯穿始终的判断}

\textbf{表示成为可优化的对象，这是唯一的结构性变化。}前面的模型族要求人先决定用什么特征，模型的假设因此是显式的、可以被质疑的；深度模型把特征交给优化器去找，于是``它假设了什么''这个问题不再有纸面上的答案，只能靠干预、迁移、探针这些间接手段去问。能力上的收获与可解释性上的损失是同一件事的两面，不是两个可以分开优化的目标。

\textbf{反向传播是链式法则，不是学习理论。}它保证的只有一件事：梯度算对了。它不保证目标函数写对了，不保证优化能到达好的区域，更不保证到达的地方能泛化。相应地，训练失败必须分三层诊断——算得对不对是数值问题，走不走得动是优化问题，学到的能不能推广是统计问题。混淆这三层，最典型的后果就是在该补数据的时候去换优化器。

\textbf{架构就是先验，而先验的强弱要匹配数据量。}卷积把局部性与平移等变硬编码进去，注意力则假设任意两个位置可以直接对话。这些都不是中性的实现细节，而是关于数据长什么样的断言。数据不足时强先验是巨大的优势，因为它把假设空间砍小了；数据充裕时它变成天花板，因为它排除了真实存在但不符合该先验的结构。这条权衡曲线上不存在普遍最优的架构，只存在与当前数据量匹配的架构。

\textbf{生成模型的失败形态由散度的方向决定。}让模型分布靠近数据分布，用哪个方向的 KL 是有后果的：一个方向重罚``该有却没有''，于是模型摊平去覆盖所有模式，代价是把质量也摊到了模式之间的空隙上，样本发糊；另一个方向重罚``不该有却有''，于是模型缩进几个模式里保证像，代价是丢掉其余。模糊与模式坍缩不是两个独立的工程故障，是同一个不对称性的两侧。

\textbf{结论要按认识论地位分级阅读。}这一部分里，等变性、链式法则、变量替换公式是定理；收敛率、残差的梯度通路是有明确条件的分析结论；而缩放律的指数、归一化为什么有效、过参数化为什么不过拟合，是被反复观察到却仍缺理论的经验规律。本部分对每一条都标注了它属于哪一类。这不是学究气——上面那张图的右端，正是外推最容易出事的地方。

\subsection{十章怎么安排}

\hyperref[ch:123]{``表示与逼近：深度和非线性究竟增加了什么''}是地基，它的主要贡献是划清通用逼近定理的边界：那个定理说的是``存在这样一个网络''，没说它能被找到、也没说需要多少神经元。\hyperref[ch:124]{``损失与反向传播：目标怎样沿计算图返回参数''}把统计目标、计算图与梯度接成一条链，其中损失函数的选择决定了网络到底在学什么这一点，比反向传播的机械过程重要得多。\hyperref[ch:125]{``优化与训练：梯度正确之后为何仍会失败''}处理的正是本部分标题里那个``仍会''——它解释了 Adam 的分母衡量的是波动而不是曲率，也解释了归一化如何把权重范数变成了一个隐蔽的学习率调节器。

接下来两章讲怎样把结构写进架构。\hyperref[ch:126]{``卷积网络：用一次交易换来的三条先验''}把局部性、平移等变、参数共享讲成一笔可以算清的交易，其中感受野是可以逐层递推算出来的，而``理论感受野不等于有效感受野''这一条解释了为什么网络名义上看得到全图却仍依赖局部纹理。\hyperref[ch:127]{``注意力与序列模型：读谁交给内容决定，位置、深度与算力另外付账''}是这一部分篇幅最重的一章：它先证明自注意力对输入置换是等变的，于是位置编码成为必然而不是技巧，再把深度稳定性与二次复杂度这两笔账逐一摊开。

再往后三章是生成模型。\hyperref[ch:128]{``变分自编码器：把不可积的后验换成可优化的下界''}给出最基本的那次替换，并说明这个下界允许一个把潜变量彻底闲置的退化解。\hyperref[ch:140]{``生成模型的统一视角：散度选在哪一边，失败就长在哪一侧''}是三章里最该完整读完的一章——它把 VAE、流、GAN、扩散放回同一个框架，读完之后前后两章都会更容易。\hyperref[ch:141]{``扩散模型：把一次难的匹配拆成很多次易的匹配''}则把那个拆解讲透，其中采样其实是在解一个常微分方程这一条，直接把它接回了数值分析。

最后两章面向交付与理解。\hyperref[ch:142]{``缩放与适配：规模的收益是幂律，代价是指数''}给出那个残酷的算术——指数在 \(0.05\) 到 \(0.1\) 这个量级，意味着损失降一半需要规模放大几个数量级；同一章里给定预算怎样在参数与数据之间分配，则是一次干净的约束优化。\hyperref[ch:129]{``AI 可解释性的数学基础：理解高维数据的流动''}放在最后，它给出的最有用的东西是一把筛子：凡是在网络的等价变换下会改变的量，描述的是坐标而不是模型。

上一部分把``用什么表示''交给人来设计，这一部分把它交给优化器。往后走，\hyperref[ch:130]{``微分方程与动力系统：从局部变化规律到整体行为''}会回过头来解释这里反复出现却没被解释的一件事——训练本身就是一个离散动力系统，学习率越过某个临界值时的那种突变，在那里有现成的语言。
